\section{变系数}\label{sec:variable-coefficients}

下面开始超出 Hodge Laplacian 的范围. 一个简单的变化, 是将一个或多个底层 Hilbert 空间 $W^k=L^2\Lambda^k$ 上的内积, 从标准 $L^2$ 内积改为加权内积:
\[
\langle\omega,\mu\rangle=\int_\Omega\langle A\omega,\mu\rangle\operatorname{vol},\qquad\omega,\mu\in L^2\Lambda^k(\Omega).
\]
这里 $A$ 是 $\Omega$ 上的函数, 对每个 $x\in\Omega$, $A(x)$ 都是 $\operatorname{Alt}^kT_x\Omega$ 上的对称正定算子 (例如, 若 $\Omega$ 是 $\mathbb R^n$ 中的区域, 则 $A$ 可视为从 $\Omega$ 到阶数为 $\binom nk$ 的对称正定矩阵空间的函数). 假定 $A$ 有上界, 且其最小特征值有正下界, 则 $W^k$ 上的加权内积与通常内积等价, 因而新的 Hilbert 复形继承了原复形的所有本质性质, 例如像是闭的. 然而, 该复形的抽象 Hodge Laplacian 有所不同. 将一般情形留给读者, 考虑 $\mathbb R^n$ 中区域上 de Rham 复形最后一个片段这一简单例子. 取上述加权空间 $L^2\Lambda^{n-1}(\Omega,A)$ 作为底层空间, 其中 $A$ 是 $n\times n$ 对称正定矩阵值函数, 复形其余部分保持不变, 因而相关片段为
\begin{equation}\label{eq:weighted-divergence-complex}
L^2\Lambda^{n-1}(\Omega,A)\xrightarrow{(\operatorname{div},H(\operatorname{div}))}L^2\Lambda^n(\Omega)\to0.
\end{equation}

这个复形的抽象 Hodge Laplace 问题的弱形式, 就是求 $\sigma\in H(\operatorname{div})$, $u\in L^2$, 使得
\[
\begin{aligned}
\langle A\sigma,\tau\rangle-\langle u,\operatorname{div}\tau\rangle&=0,&\tau&\in H(\operatorname{div}),\\
\langle\operatorname{div}\sigma,v\rangle&=\langle f,v\rangle,&v&\in L^2.
\end{aligned}
\]
这是下述问题相应的标准弱形式
\[
A\sigma+\operatorname{grad}u=0,\qquad\operatorname{div}\sigma=f,
\]
或者等价地, $-\operatorname{div}C\operatorname{grad}u=f$, $\sigma=-C\operatorname{grad}u$, 其中 $C=A^{-1}$, 并附有 $\partial\Omega$ 上的 Dirichlet 边界条件 $u=0$; 在此弱形式中, 它是自然边界条件. 简言之, 这一变系数偏微分算子仍是某个 Hilbert 复形的抽象 Hodge Laplacian, 尽管这个 Hilbert 复形与 $L^2$ de Rham 复形稍有不同. 可以用 $A$ 为恒等映射时的相同有限元空间进行离散化, 分析无需额外考察即可推广. 其他许多变系数椭圆型偏微分方程也可用类似方式处理.
